\section{Three Things That Only Break Together}
\label{sec:ch08-only-breaks-together}
\index{composition!failures a subgraph cannot see|(}%

Both subgraphs work. Each publishes a schema, answers its own root fields,
resolves representations and passes its tests. Hand the two schemas to a
composer and none of that is enough.

Composition is chapter~\ref{ch:composition}'s subject and I am not going to
spend it here. But three of the errors I met belong to this chapter rather than
to that one. Each is a setting on a subgraph, each is on by default, and no
test of a single service can reach any of them. Until a schema has been put
beside another one, nobody has really read it.

\subsection*{Two services cannot both have \code{node}}

\index{Node interface@\code{Node} interface!and federation}%
The first error names a field neither service wrote:

\begin{minted}[fontsize=\footnotesize,breaklines]{text}
The Object "Query" defines the same fields in multiple subgraphs without the
"@shareable" directive:
 The field "node" is defined in the following subgraphs: "catalog", "mosaic".
 However, it is not declared "@shareable" in any of them.
 The field "nodes" is defined in the following subgraphs: "catalog", "mosaic".
 However, it is not declared "@shareable" in any of them.
\end{minted}

Chapter~\ref{ch:schema-design-that-survives-change} gave Mosaic the \code{Node}
interface and the two root fields that come with it. Both services still have
types that implement it, so both ask for global object identification, so both
declare \code{node} and \code{nodes}. Federation~v2 refuses a field that two
subgraphs define unless one of them says it is
shareable~\autocite{apollo2026compositionrules}.

\index{shareable@\code{@shareable}!when it would be a lie}%
The tempting fix is to mark them shareable and move on. It would compose, and it
would be a lie. \code{@shareable} means two subgraphs can resolve the field
interchangeably, and these two cannot: Catalog's \code{node} knows how to fetch
a \code{Product} and has never heard of an \code{Order}, and Mosaic's is the
mirror image. A router free to send \code{node(id:)} to whichever service was
convenient would answer null for roughly half the identifiers it was given, and
which half would depend on the query plan.

So the field goes. One argument to one call, in both services:

\begin{minted}{csharp}
.AddGlobalObjectIdentification(registerNodeInterface: false)
\end{minted}

That reads like it removes more than it does, and I checked exactly what
survives. The node identifier serialiser stays registered, so every \code{[ID]}
still encodes and the federation keys of
section~\ref{sec:ch08-the-key-you-already-designed} are unaffected. The
\code{Node} interface stays in both schemas. \code{Customer}, \code{Order} and
\code{Review} still say \code{implements Node}, and so does Catalog's
\code{Product}. The node resolvers still compile and are still there. What goes
is \code{Query.node} and \code{Query.nodes}, and nothing else.

\index{Node interface@\code{Node} interface!what the split costs}%
That is a real loss and it should be recorded as one rather than smoothed over.
Global object identification was two promises: that every object has an
identifier you can hold on to, and that there is one field you can hand it back
to. The first survives the split intact. The second does not, and it does not
survive because it was never a property of a schema so much as of a process that
could see everything. Chapter~\ref{ch:hard-modeling-problems} is where a
federated graph gets a \code{node} field back, and it is a harder problem than
it looks, because whoever answers it has to know which subgraph owns a type it
has only ever seen as base64.

\subsection*{The cost directives are not the ones the composer knows}

\index{cost analysis!and composition}%
The second error arrives once per field, for most of the fields in both
services:

\begin{minted}[fontsize=\footnotesize,breaklines]{text}
The value ""10"" provided to argument "@cost(weight: ...)" is not a valid "Int!" type.
The definition for "@listSize" does not define the following argument that is
provided: "slicingArgumentDefaultValue"
\end{minted}

HotChocolate stamps \code{@cost(weight: "10")} on every resolver-backed field by
default, with the weight as a string. The composer in wgc 0.129.7 carries a
definition whose weight is an \code{Int!}. I am not going to adjudicate which of
them is reading the cost specification correctly, because it does not matter to
anybody trying to ship: two tools disagree about the type of one directive
argument, and the first thing in the world that ever compares them is the
composer.

The \code{@listSize} half is not even a disagreement. HotChocolate's own option
for that argument is called \code{ApplySlicingArgumentDefaultValue} and its
documentation comment reads \enquote{Defines if the non-spec slicing argument
default value shall be applied}. The library knows it is emitting something
outside the specification, has a switch for it, and leaves the switch on.

Both switches live in one call, and both services need it:

\begin{minted}{csharp}
.ModifyCostOptions(options =>
{
    options.ApplyCostDefaults = false;
    options.ApplySlicingArgumentDefaultValue = false;
})
\end{minted}

\index{cost analysis!what turning defaults off costs}%
What that turns off is narrower than the name suggests, and I checked this too
because I did not want to disable a security control by accident.
\code{ApplyCostDefaults} and \code{SkipAnalyzer} are separate settings on the
same options class. The analyzer is still in the pipeline, the limits are still
enforced, and an explicit \code{[Cost]} still counts. What goes away is the
automatic weight on every resolver-backed field, and with it the cost numbers
chapter~\ref{ch:schema-design-that-survives-change} printed for
\code{Product.reviews}. Those measurements stand for the schema that produced
them and do not survive this setting, which is worth saying plainly rather than
letting a reader find out by rerunning them.
Chapter~\ref{ch:security-hardening} is where cost is a subject rather than a
composition problem, and it will have to set weights deliberately.

\subsection*{The type nobody wrote}

\index{PageCursor@\code{PageCursor}!and shareability}%
The third error names a type that appears in neither codebase:

\begin{minted}[fontsize=\footnotesize,breaklines]{text}
The Object "PageCursor" defines the same fields in multiple subgraphs without the
"@shareable" directive:
 The field "page" is defined in the following subgraphs: "catalog", "mosaic".
 The field "cursor" is defined in the following subgraphs: "catalog", "mosaic".
\end{minted}

Both services have a connection, so both have a \code{PageInfo}, so both have
the \code{PageCursor} that \code{PageInfo.forwardCursors} returns. Two subgraphs
declaring the same object is the same problem as \code{node}, and this time
\code{@shareable} is the right answer: a page cursor is a value with no identity
and either service can produce one.

What makes it worth a page is that HotChocolate already knows about this
problem and fixes half of it. Read \code{FederationTypeInterceptor} at tag
\code{16.6.0} and the comment is right there, typo and all:

\begin{minted}[fontsize=\footnotesize]{csharp}
// if we find a PagingInfo we will make all fields sharable.
if (configuration is ObjectTypeConfiguration typeCfg
    && typeCfg.Name.Equals(PageInfoType.Names.PageInfo))
\end{minted}

\index{HotChocolate!PageInfo auto-shareable@\code{PageInfo} auto-shareable}%
It matches one type by name. \code{PageInfo} comes out marked and every one of
its fields is shareable in the published schema, which is why that type is not
in the error. \code{PageCursor} arrives in the schema through \code{PageInfo},
gets missed, and takes the blame. Any two HotChocolate subgraphs that both have
a connection will hit this.

The fix is two attributes on an empty class, in both services, and it is the
same mechanism the rest of this repository already uses to add something to a
type it does not own:

\begin{minted}{csharp}
[Shareable]
[ObjectType<PageCursor>]
public static partial class PageCursorExtension;
\end{minted}

\code{@shareable} on an object makes every field on it shareable, so there is
nothing else to write. The published schema then says
\code{type PageCursor @shareable} and the composer stops caring.

\subsection*{What that leaves}

With those three settled the two schemas compose, and I ran the composer once to
be sure of it before writing any of this down. That is the whole of what
chapter~\ref{ch:the-first-cut-extracting-catalog} does with composition: one
command, no output shown, used as a gate rather than as a subject. The
repository's verification script runs it on every pass for the same reason. A
pair of subgraphs that has quietly stopped composing is precisely the breakage
that no test of either service can find.

The pattern across all three is the thing to carry forward. None of them was a
mistake in the extraction. Each was a default that is correct for one service
and wrong for two, and each became visible only when a second schema existed to
compare against.

So the expensive part of the first cut was not moving the code. It was finding
out that a schema I had been shipping for seven chapters had opinions in it that
only mattered once somebody else had opinions too.
Chapter~\ref{ch:composition} is where those get read properly.

\index{composition!failures a subgraph cannot see|)}%
